Os maiores percentuais de rendimentos de outras fontes foram observados em Macambira, Pedrinhas, Gararu, conforme a Tabela~\ref{tab:item5-outras-fontes}. A correlação de Pearson entre o rendimento total médio e a participação dessas fontes, apresentada na Figura~\ref{fig:item5-total-participacao}, foi $r = -0{,}468$ ($p < 0{,}001$). A associação negativa indica que municípios com maior rendimento total tendem a apresentar menor participação relativa de outras fontes; portanto, essa participação é relativamente mais elevada em municípios de menor rendimento. Isso não demonstra, por si só, dependência da economia municipal em relação ao governo: P1100 reúne outras fontes de renda sem permitir separar aposentadorias, pensões, PBF ou demais componentes, e a correlação municipal não estabelece causalidade.
